\documentclass[11pt]{article}
\usepackage[margin=1in]{geometry}
\usepackage{amsmath,amssymb,amsthm,hyperref}
\newtheorem{theorem}{Theorem}
\newtheorem{proposition}[theorem]{Proposition}
\title{Common positive zeros of the derivative-polynomial combinations}
\author{Principia Math In Progress}
\date{11 September 2026}
\begin{document}
\maketitle
\begin{abstract}
Conjecture 2.8 of Erd\'elyi, Rosenblatt, and Rosenblatt is false already
at derivative order two. With coefficients $(18,30,25)$ the two polynomials
have the common positive zero $1/2$. More generally, every prescribed
positive point is a common zero for suitable strictly positive coefficients.
The counterexample is minimal in derivative order.
\end{abstract}

\section*{Task 1: Resolve MathDB problem 355875}
In the notation of \cite{source}, define
\[
 A_0(x)=\frac{x}{(1+x^2)^{3/2}},\qquad
 B_0(x)=\frac1{(1+x^2)^{3/2}},
\]
and write the ordinary derivatives as
\[
 A_0^{(l)}(x)=\frac{P_l(x)}{(1+x^2)^{(2l+3)/2}},\qquad
 B_0^{(l)}(x)=\frac{Q_l(x)}{(1+x^2)^{(2l+3)/2}}.
\]
Given real coefficients $\gamma_0,\ldots,\gamma_L$, not all zero, put
\begin{align*}
 F_\gamma(x)&=\sum_{l=0}^L
 \gamma_l(1+x^2)^{L-l}x^lP_l(x),\\
 G_\gamma(x)&=\sum_{l=0}^L
 \gamma_l(1+x^2)^{L-l}x^lQ_l(x).
\end{align*}
Conjecture 2.8, page 11, asserts that their positive real zero sets are
disjoint. The surrounding discussion of common zeros and interlacing
specifies this meaning of distinctness.

\begin{theorem}\label{thm:example}
For $L=2$ and $(\gamma_0,\gamma_1,\gamma_2)=(18,30,25)$,
\begin{align}
 F_\gamma(x)&=3x(4x^2-1)(9x^2-16),\label{eq:F}\\
 G_\gamma(x)&=3(4x^2-1)(19x^2-6).\label{eq:G}
\end{align}
Both polynomials are nonzero and have a simple zero at $x=1/2$.
Thus the conjecture fails with strictly positive coefficients and a common
zero strictly between zero and one.
\end{theorem}

\begin{proof}
Direct differentiation gives the following table:
\[
\begin{array}{c|cc}
 l&P_l(x)&Q_l(x)\\\hline
 0&x&1\\
 1&1-2x^2&-3x\\
 2&6x^3-9x&12x^2-3.
\end{array}
\]
For example, each numerator sequence satisfies
\[
 R_{l+1}(x)=(1+x^2)R_l'(x)-(2l+3)xR_l(x),
\]
with its respective initial value. Substitution of the table and the
three coefficients yields
\begin{align*}
 F_\gamma(x)
 &=18x(1+x^2)^2+30x(1+x^2)(1-2x^2)
       +25x^2(6x^3-9x)\\
 &=48x-219x^3+108x^5
 =3x(4x^2-1)(9x^2-16),\\
 G_\gamma(x)
 &=18(1+x^2)^2-90x^2(1+x^2)+25x^2(12x^2-3)\\
 &=18-129x^2+228x^4
 =3(4x^2-1)(19x^2-6).
\end{align*}
The factor $4x^2-1$ vanishes simply at $1/2$. At that point the remaining
factors are nonzero: $x=1/2$, $9x^2-16=-55/4$, and
$19x^2-6=-5/4$. This proves the theorem.
\end{proof}

\begin{proposition}\label{prop:family}
For every $r>0$, let $u=r^2$ and take $L=2$ with
\[
 (\gamma_0,\gamma_1,\gamma_2)
 =\bigl(3u(2u+1),\ 6u(1+u),\ (1+u)^2\bigr).
\]
All three coefficients are strictly positive, and
\begin{align*}
 F_\gamma(x)&=3x(x^2-u)\bigl((u+2)x^2-4u-3\bigr),\\
 G_\gamma(x)&=3(x^2-u)\bigl((3u+4)x^2-2u-1\bigr).
\end{align*}
In particular, $r$ is a common zero. Conversely, for $L=2$, all nonzero
coefficient vectors giving a common zero at $r$ are scalar multiples of
this vector. There is no counterexample at $L=0$ or $L=1$.
\end{proposition}

\begin{proof}
The displayed factorizations follow by substitution in the derivative
table. For completeness, at a common zero $r>0$ the equations obtained
by dividing $F_\gamma(r)$ by $r$ and setting $G_\gamma(r)=0$ are
\begin{align*}
 \gamma_0(1+u)^2+\gamma_1(1+u)(1-2u)
       +\gamma_2u(6u-9)&=0,\\
 \gamma_0(1+u)^2-3\gamma_1u(1+u)
       +\gamma_2u(12u-3)&=0.
\end{align*}
Their difference gives
$\gamma_1(1+u)^2=6\gamma_2u(1+u)$.
Substitution in the second equation then gives
$\gamma_0(1+u)^2=3\gamma_2u(2u+1)$.
If $\gamma_2=0$, all coefficients vanish, which is forbidden.
These relations prove the classification.

For $L=0$, the pair is $(\gamma_0x,\gamma_0)$ and has no common
positive zero. For $L=1$, write
\[
 F_\gamma(x)/x=\gamma_0(1+x^2)+\gamma_1(1-2x^2),\qquad
 G_\gamma(x)=\gamma_0(1+x^2)-3\gamma_1x^2.
\]
At a common positive zero, subtraction gives
$\gamma_1(1+x^2)=0$, hence $\gamma_1=0$ and then $\gamma_0=0$.
This again contradicts the hypothesis.
\end{proof}

\paragraph{Scope and evidence.}
The explicit conjecture permits arbitrary real coefficients. Our example
also satisfies strict positivity, and the common zero is neither the origin
nor the value one. Both individual zeros at $1/2$ are simple; the failure is
that the two zero sets intersect. No conclusion about the finiteness of
physical electrical-field equilibria follows from this polynomial
counterexample alone.

The same-agent proof audit, exact symbolic verifier, and dated literature
audit accompany this file. The current arXiv and author-hosted manuscripts
retain Conjecture 2.8. MathDB was checked open with no posted claims on
11 September 2026, and no existing resolution was found in the documented
searches. No external peer review is claimed.

\section*{Task 2: Relate the coefficients to point charges}
Determine whether the positive coefficient examples above can arise from
the source's leading moment data for a finite point-charge configuration;
for realizable examples, determine whether actual zero branches have a
common asymptotic direction.

\section*{Task 3: Characterize coefficients with disjoint zero sets}
For each $L\geq3$, characterize the coefficient vectors for which
$F_\gamma$ and $G_\gamma$ have no common positive real zero, and identify
useful explicit conditions ensuring that property.

\section*{Task 4: Bound the number of common positive zeros}
Determine the largest possible number of distinct common positive zeros
for nonzero $F_\gamma,G_\gamma$ of order $L$, and classify coefficient
vectors attaining that bound.

\section*{Task 5: Analyze perturbations of the common zeros}
For coefficient vectors near the family in Proposition~\ref{prop:family},
derive the first-order splitting of the two roots near $r$ and obtain
uniform separation estimates away from parameters producing multiple zeros.

\begin{thebibliography}{9}
\bibitem{source} T. Erd\'elyi, J. Rosenblatt, and R. Rosenblatt,
\emph{Asymptotic Directions for the Zero Sets of the Components of an
Electrical Field from a Finite Number of Point Charges on the Plane Part II},
\href{https://arxiv.org/abs/2208.12857}{arXiv:2208.12857v1}, 26 August 2022,
Section 2.7, Conjecture 2.8, p.~11.
\end{thebibliography}
\end{document}
